\documentclass[12pt]{article}
\usepackage[utf8]{inputenc}
\usepackage[T1]{fontenc}
\usepackage{hyperref}
\title{Panawa Language Resources for K-12 Curriculum Development}
\author{Compiled by Claude Sonnet 5.5 (claude-sonnet-5-5) for Andr\'as Kornai}
\date{}
\begin{document}
\maketitle

\section{Introduction}
Panawa (ISO 639-3 pwb; also Bujiyel, Bujiye, Bugel, iPanawa) is an East Kainji language of the Shammo cluster, spoken in Toro LGA, Bauchi State, Nigeria. Its closest relative is reported to be Boze. Speaker estimates vary widely between sources (roughly 3,000 to 6,100). Web searches found very little material specific to the language: no primer, textbook, dictionary or school programme was found.

\section{Primary Educational Resources}
No primers, textbooks or school materials in Panawa were found in search results. Hausa is widely used in the area and some English is spoken, so any curriculum would currently have to be built from scratch.

\section{Linguistic and Reference Materials}
\begin{itemize}
\item Wikipedia, ``Panawa language'': overview of classification, location and speaker numbers. \url{https://en.wikipedia.org/wiki/Panawa_language}
\item Glottolog bibliographic record for Blench and Nengel (2004), a reference on the language. Only the catalogue record was seen, not the text itself. \url{https://glottolog.org/resource/reference/id/18723}
\end{itemize}

\section{Religious and Community Resources}
\begin{itemize}
\item Joshua Project, Panawa people group in Nigeria: population, religion (mainly Sunni Muslim) and scripture status (translation started). \url{https://joshuaproject.net/people_groups/19670}
\item Joshua Project language page for pwb. \url{https://joshuaproject.net/languages/pwb}
\end{itemize}

\section{Digital and Audio Resources}
\begin{itemize}
\item Global Recordings Network page for Panawa, which lists audio materials in the language. \url{https://globalrecordings.net/en/language/pwb}
\end{itemize}

\section{Curriculum Development Context}
Panawa is not known to be taught in schools, and published documentation appears sparse. A curriculum would need new fieldwork with the community, an agreed orthography, and the Global Recordings audio as a possible starting point. The Blench and Nengel work should be located and checked.
\end{document}